\documentclass[11pt]{article}
\usepackage[margin=1in]{geometry}
\usepackage{amsmath,amssymb,amsthm}
\usepackage{enumitem}
\usepackage[hidelinks]{hyperref}
\usepackage{microtype}
\usepackage{lmodern}

\newtheorem{theorem}{Theorem}
\newtheorem{lemma}{Lemma}
\newcommand{\Ls}{L_s}
\newcommand{\lmax}{\lambda_{\max}}
\newcommand{\tmax}{T_{\max}}
\newcommand{\Tau}{\mathcal{T}}

\title{A Proof of the Source-Corrected\\Written on the Wall II Graph Conjecture 160}
\author{AI-assisted proof search and write-up}
\date{July 23, 2026}

\begin{document}
\maketitle

\begin{abstract}
For a finite nontrivial connected simple graph $G$, let $\Ls(G)$ be the
maximum number of leaves in a spanning tree, let
$\lambda(v)=\alpha(G[N(v)])$, let $T(v)$ be the number of triangles
containing $v$, and let $\chi_{C_4}(G)$ be $1$ when $G$ contains no
(not necessarily induced) $4$-cycle and $0$ otherwise.  We prove the
source-corrected Written on the Wall II Conjecture 160:
\[
  \Ls(G)\ge \max_v\lambda(v)+\chi_{C_4}(G)\max_v T(v).
\]
The proof is constructive.  In the $C_4$-free case, every open
neighborhood is a matching together with isolated vertices.  Two
suitably chosen stars, joined by a geodesic, then produce a tree
subgraph with the required number of leaves; a short exceptional case
is handled by rerouting through a degree-two triangle.
\end{abstract}

\section{Statement and notation}

Let $G$ be a finite nontrivial connected simple graph.  Write
\begin{itemize}[leftmargin=2em]
  \item $\Ls(G)$ for the maximum number of leaves among spanning trees of $G$;
  \item $\lambda(v)=\alpha(G[N(v)])$, the independence number of the graph
        induced by the open neighborhood of $v$;
  \item $T(v)$ for the number of triangles containing $v$;
  \item $\lmax=\max_v\lambda(v)$ and $\tmax=\max_vT(v)$;
  \item $\chi_{C_4}(G)=1$ if $G$ contains no $4$-cycle as a subgraph,
        and $\chi_{C_4}(G)=0$ otherwise.
\end{itemize}
The corrected historical statement is
\begin{equation}\label{eq:main}
  \Ls(G)\ge \lmax+\tmax\chi_{C_4}(G).
\end{equation}
The $4$-cycle in this definition need not be induced.

\section{Two elementary lemmas}

\begin{lemma}[Subtree extension]\label{lem:extend}
If $H$ is a tree subgraph of a connected graph $G$, then $H$ is
contained in a spanning tree of $G$ having at least as many leaves as $H$.
\end{lemma}

\begin{proof}
While the current tree is not spanning, connectivity gives an edge
$uv$ with $u$ in the current tree and $v$ outside it.  Add $v$ and the
edge $uv$.  If $u$ was not a leaf, the operation creates one new leaf.
If $u$ was a leaf, $u$ ceases to be a leaf but $v$ becomes one, so the
number of leaves does not decrease.  Iteration gives a spanning tree.
\end{proof}

Thus it is enough to construct any tree subgraph with the desired
number of leaves.

\begin{lemma}[Neighborhood structure]\label{lem:local}
If $G$ contains no $4$-cycle, then $G[N(v)]$ has maximum degree at most
one for every vertex $v$.
\end{lemma}

\begin{proof}
If some $a\in N(v)$ had two distinct neighbors $b,c\in N(v)$, then
$b-a-c-v-b$ would be a $4$-cycle.
\end{proof}

Consequently, in a $C_4$-free graph, $G[N(v)]$ is a disjoint union of
edges and isolated vertices.  Each edge of $G[N(v)]$ is a triangle
through $v$.  If $G[N(v)]$ contains $T(v)$ edges and $s(v)$ isolated
vertices, then
\[
  d(v)=2T(v)+s(v),\qquad \lambda(v)=T(v)+s(v).
\]
Hence
\begin{equation}\label{eq:degree}
  d(v)=\lambda(v)+T(v),
\end{equation}
and
\begin{equation}\label{eq:lambdaT}
  \lambda(v)\ge T(v).
\end{equation}
Also, two vertices in a $C_4$-free graph have at most one common
neighbor.

\section{Tree subgraphs built from two stars}

Let $x$ and $y$ be distinct vertices in a connected $C_4$-free graph.
The following constructions will be used repeatedly.

\begin{enumerate}[label=\textbf{(\roman*)}, leftmargin=2.8em]
\item If $d(x,y)\ge2$, take a shortest $x$--$y$ path, all edges from
$x$ to $N(x)$, and all edges from $y$ to $N(y)$.  Geodesicity prevents
an internal path vertex, other than the first or last one, from being a
neighbor of an endpoint.  If $d(x,y)=2$, the middle vertex is the only
common neighbor of $x$ and $y$, since a second common neighbor would
create a $4$-cycle.  The selected edges therefore form a tree with at
least
\begin{equation}\label{eq:far}
  d(x)+d(y)-2
\end{equation}
leaves.

\item If $x$ and $y$ are adjacent and have no common neighbor, their
full stars form a double-star with at least
\begin{equation}\label{eq:adj0}
  d(x)+d(y)-2
\end{equation}
leaves.

\item If $x$ and $y$ are adjacent and have a common neighbor $z$, then
$z$ is unique.  Take both full stars and omit one of the two selected
edges joining $z$ to $x,y$.  The resulting tree has at least
\begin{equation}\label{eq:adj1}
  d(x)+d(y)-3
\end{equation}
leaves.
\end{enumerate}

By Lemma~\ref{lem:extend}, each of these lower bounds also applies to
$\Ls(G)$.

\section{Proof of the conjecture}

\begin{theorem}
Every finite nontrivial connected simple graph satisfies
Equation~\eqref{eq:main}.
\end{theorem}

\begin{proof}
Choose vertices $x,y$ such that
\[
  \lambda(x)=\lmax=: \Lambda,
  \qquad
  T(y)=\tmax=: \Tau.
\]

First suppose that $G$ contains a $4$-cycle.  Then
$\chi_{C_4}(G)=0$.  A maximum independent set in $G[N(x)]$, joined to
$x$, gives a star with at least $\Lambda$ leaves.  Lemma~\ref{lem:extend}
therefore gives $\Ls(G)\ge\Lambda$, which proves the desired bound in
this case.

Now suppose that $G$ is $C_4$-free.  Put
\[
  t_x=T(x),\qquad \lambda_y=\lambda(y).
\]
By Equation~\eqref{eq:degree},
\begin{equation}\label{eq:extremaldegrees}
  d(x)=\Lambda+t_x,
  \qquad
  d(y)=\lambda_y+\Tau.
\end{equation}
If $\Tau=0$, the full star at $x$ has at least
$d(x)=\Lambda+t_x\ge\Lambda$ leaves, so assume $\Tau\ge1$.

If $x=y$, then $t_x=\Tau$, and the full star at $x$ has
$d(x)=\Lambda+\Tau$ leaves.  Hence assume $x\ne y$.

\medskip
\noindent\textbf{Case 1: $d(x,y)\ge2$.}
If $t_x+\lambda_y\ge2$, then Equations~\eqref{eq:far} and
\eqref{eq:extremaldegrees} give
\[
\begin{aligned}
  \Ls(G)
  &\ge d(x)+d(y)-2\\
  &=\Lambda+\Tau+(t_x+\lambda_y-2)\\
  &\ge\Lambda+\Tau.
\end{aligned}
\]
Suppose instead that $t_x+\lambda_y<2$.  Since $\Tau\ge1$ and
Equation~\eqref{eq:lambdaT} gives $\lambda_y\ge\Tau$, necessarily
\[
  t_x=0,\qquad \lambda_y=\Tau=1.
\]
Lemma~\ref{lem:local} then implies that $d(y)=2$, and the two neighbors
$p,z$ of $y$ are adjacent.  Choose a shortest path from $x$ to $y$ so
that $p$ is the predecessor of $y$.  Take the full star at $x$, this
path, and the edge $pz$, while omitting $yz$.  This is a tree.  Its
leaves include the $d(x)-1$ neighbors of $x$ not used to continue the
path, as well as both $y$ and $z$.  It therefore has
\[
  d(x)-1+2=d(x)+1=\Lambda+1=\Lambda+\Tau
\]
leaves.

\medskip
\noindent\textbf{Case 2: $x$ and $y$ are adjacent and have no common
neighbor.}
Here $x$ is an isolated vertex of $G[N(y)]$.  The same neighborhood
also contains the $\Tau\ge1$ independent matching edges corresponding
to the triangles through $y$.  Thus
\[
  \lambda_y\ge\Tau+1\ge2.
\]
In particular $t_x+\lambda_y\ge2$, and Equations~\eqref{eq:adj0} and
\eqref{eq:extremaldegrees} yield $\Ls(G)\ge\Lambda+\Tau$.

\medskip
\noindent\textbf{Case 3: $x$ and $y$ are adjacent and have a common
neighbor.}
The common neighbor is unique.  The triangle through $x,y$ gives
$t_x\ge1$, while maximality of $\Tau$ gives $t_x\le\Tau$.

If $t_x=\Tau$, the full star at $x$ has
\[
  d(x)=\Lambda+t_x=\Lambda+\Tau
\]
leaves.  If $t_x<\Tau$, then $\Tau\ge2$.  Equation~\eqref{eq:lambdaT}
gives $\lambda_y\ge\Tau\ge2$, and therefore
$t_x+\lambda_y\ge3$.  Equations~\eqref{eq:adj1} and
\eqref{eq:extremaldegrees} now give
\[
\begin{aligned}
  \Ls(G)
  &\ge d(x)+d(y)-3\\
  &=\Lambda+\Tau+(t_x+\lambda_y-3)\\
  &\ge\Lambda+\Tau.
\end{aligned}
\]
This completes every case.
\end{proof}

\section{Computational audit}

The accompanying program \texttt{verify\_wowii160.py} implements the
constructive proof.  For an input graph it computes the extremal
vertices $x,y$, constructs the appropriate tree subgraph in the case
analysis above, extends it to a spanning tree without losing leaves,
and checks the claimed inequality.  It also exhaustively checks all
995 connected unlabeled graphs of orders two through seven contained
in the NetworkX graph atlas, followed by a deterministic random test.
The computation is an audit of the construction, not a substitute for
the proof.

\section*{Status note}

This note proves the corrected historical statement in which the
$C_4$ term is the binary $C_4$-free characteristic.  A separate source
audit found that an earlier formalization used the number of induced
$4$-cycles instead; that is not the statement proved here.  A search of
the current Formal Conjectures pull requests found a correction of the
statement but no submitted proof of the corrected conjecture at the
time of this write-up.  Mathematical priority should nevertheless be
assessed independently by specialists and through a full literature
search.

\end{document}
